\documentclass[11pt]{article}
\usepackage{../../latex/stat4ml_style}

\title{Chapter 0: Matrix Algebra Review \\ \large Part I --- Statistics Foundation for ML}
\author{Omid Safarzadeh}
\date{}

\begin{document}
\maketitle

\begin{abstract}
\noindent
This is chapter zero on purpose: everything from Chapter 2 onward (moment
matrices, covariance matrices, multivariate distributions, and eventually
regression coefficients and neural-network weight matrices) is written in
matrix notation, and none of it is readable without this chapter first.
This chapter reviews matrix multiplication, determinants, inverses,
transposes, partitioned matrices, eigenvalues/eigenvectors, the four
standard matrix decompositions (LU, Cholesky, QR, SVD), and matrix
differentiation -- the exact toolkit \S4 (Matrix Notation for Moments) of
Chapter 2 assumes you already have.

\emph{Acknowledgement: this chapter expands the STAT4ML prerequisite
slides (``Matrix Calculus''), which drew on Barnes (2009) and Wikipedia's
matrix decomposition article; see the reference list at the end.}
\end{abstract}

\tableofcontents

\section{Matrix Multiplication}

\begin{definition}[Matrix product]
Let $A$ be $m \times n$ and $B$ be $n \times p$. The product $C = AB$ is
the $m \times p$ matrix with entries
\[
  c_{ij} = \sum_{k=1}^n a_{ik} b_{kj}, \qquad i = 1,\dots,m,\ j = 1,\dots,p.
\]
\end{definition}

Matrix multiplication is only defined when the inner dimensions match ($A$
has as many columns as $B$ has rows); the result takes the outer
dimensions. This single rule is why a weight matrix in a neural network
must have shape (output features $\times$ input features) -- so it can
left-multiply an input vector of the right length.

\section{Determinants}

For a $2\times 2$ matrix $A = \begin{psmallmatrix} a & b \\ c & d
\end{psmallmatrix}$, $\det A = ad - bc$.

\begin{theorem}[Properties of determinants]
For square matrices $A, B$ of the same size:
\begin{enumerate}
  \item $\det I = 1$.
  \item Exchanging two rows of a matrix flips the sign of its determinant.
  \item The determinant is a linear function of each row separately: for
        $2\times 2$, $\begin{vmatrix} ta+a' & tb+b' \\ c & d\end{vmatrix}
        = t\begin{vmatrix}a&b\\c&d\end{vmatrix} +
        \begin{vmatrix}a'&b'\\c&d\end{vmatrix}$.
  \item If two rows of $A$ are equal, $\det A = 0$.
  \item Subtracting $t$ times row $i$ from row $j$ ($i \ne j$) does not
        change the determinant.
  \item If $A$ has an all-zero row, $\det A = 0$.
  \item The determinant of a triangular matrix is the product of its
        diagonal entries.
  \item $\det(AB) = (\det A)(\det B)$.
  \item $\det(A^\top) = \det A$.
\end{enumerate}
\end{theorem}

\begin{tryit}
Property 2 above is exactly why Gaussian elimination with row swaps has
to track a sign flip when computing a determinant via row reduction --
implement \texttt{det2x2} in this chapter's exercises and confirm
property 8 numerically on two random $2\times 2$ matrices.
\end{tryit}

\section{The Inverse Matrix}

\begin{definition}[Invertible matrix]
An $n\times n$ matrix $A$ is \emph{invertible} if there exists an $n
\times n$ matrix $B$ with $AB = BA = I_n$. When it exists, $B$ is unique,
called the \emph{inverse} of $A$, and denoted $A^{-1}$.
\end{definition}

For a $2\times 2$ matrix, $A^{-1} = \dfrac{1}{\det A}
\begin{psmallmatrix} d & -b \\ -c & a \end{psmallmatrix}$, which only
exists when $\det A \ne 0$ -- a matrix with zero determinant is called
\emph{singular} and has no inverse. This is the same condition that shows
up later as ``$X^\top X$ must be invertible'' for ordinary least squares
to have a unique closed-form solution (Part II, Chapter 2).

\section{Matrix Transpose}

\begin{definition}[Transpose]
The transpose of $A$, written $A^\top$, swaps $A$'s row and column
indices: $(A^\top)_{ij} = A_{ji}$.
\end{definition}

\begin{example}
For $A = \begin{psmallmatrix}1&2\\3&4\\5&6\end{psmallmatrix}$,
$A^\top = \begin{psmallmatrix}1&3&5\\2&4&6\end{psmallmatrix}$.
\end{example}

\begin{theorem}[Properties of transpose and inverse]
\begin{enumerate}
  \item $(A^\top)^\top = A$.
  \item $(A \pm B)^\top = A^\top \pm B^\top$.
  \item $(\kappa A)^\top = \kappa A^\top$ for a scalar $\kappa$.
  \item $(AB)^\top = B^\top A^\top$ (order reverses!).
  \item If $C = AB$ with $A, B$ both invertible, $C^{-1} = B^{-1}A^{-1}$
        (order reverses here too).
\end{enumerate}
\end{theorem}

\section{Partitioned Matrices}

\begin{definition}
Write $A$ ($m\times n$) in block form
$A = \begin{psmallmatrix} B & C \\ D & E \end{psmallmatrix}$,
where $B$ is $m_1\times n_1$, $E$ is $m_2\times n_2$, $C$ is $m_1\times
n_2$, $D$ is $m_2\times n_1$, $m_1+m_2=m$, $n_1+n_2=n$. This is a
\emph{partition} of $A$.
\end{definition}

Partitioned matrices are exactly the shape of the block covariance matrix
$\mathrm{Cov}\begin{psmallmatrix}Y\\Z\end{psmallmatrix} =
\begin{psmallmatrix}\mathrm{Cov}(Y) & \mathrm{Cov}(Y,Z) \\
\mathrm{Cov}(Z,Y) & \mathrm{Cov}(Z)\end{psmallmatrix}$ from Chapter 2,
\S4 -- this chapter's block-matrix notation and that chapter's block
covariance notation are literally the same object.

\section{Eigenvalues and Eigenvectors}

Almost every vector changes direction when multiplied by a matrix $A$.
The exceptional vectors that do not -- that come back out parallel to
themselves -- are \keyterm{eigenvectors}.

\begin{definition}[Eigenvalue, eigenvector]
For a square matrix $A$, a nonzero vector $x$ is an \emph{eigenvector} of
$A$ with \emph{eigenvalue} $\lambda$ if
\[
  Ax = \lambda x.
\]
$\lambda$ tells you whether $x$ is stretched, shrunk, reversed, or left
unchanged by $A$.
\end{definition}

\begin{example}
Let $A = \begin{psmallmatrix}0.8 & 0.3 \\ 0.2 & 0.7\end{psmallmatrix}$.
Solving $\det(A - \lambda I) = 0$:
\[
  \det\begin{psmallmatrix}0.8-\lambda & 0.3 \\ 0.2 & 0.7-\lambda
  \end{psmallmatrix} = \lambda^2 - \tfrac{3}{2}\lambda + \tfrac{1}{2}
  = (\lambda - 1)(\lambda - \tfrac{1}{2}) = 0,
\]
so $\lambda_1 = 1$, $\lambda_2 = 1/2$.
\end{example}

\begin{theorem}[Trace and determinant, $2\times 2$]
For a $2\times 2$ matrix $A$ with eigenvalues $\lambda_1, \lambda_2$
(possibly complex, possibly repeated):
\[
  \lambda_1 + \lambda_2 = \operatorname{trace}(A) = a_{11}+a_{22},
  \qquad
  \lambda_1 \lambda_2 = \det A.
\]
\end{theorem}
This gives a fast sanity check for any eigenvalue computation you write:
the eigenvalues must sum to the trace and multiply to the determinant --
Exercise 3 in this chapter builds exactly this check into its test suite.

\section{Matrix Decompositions}

A \keyterm{matrix decomposition} factors a matrix into a product of
(usually simpler, structured) matrices. Decompositions related to solving
linear systems: LU, Cholesky, QR. A decomposition based on eigenvalues:
SVD.

\subsection{LU Decomposition}

\begin{definition}
The \emph{LU decomposition} writes a square matrix as $A = LU$, with $L$
lower triangular and $U$ upper triangular.
\end{definition}

\begin{example}
For $A = \begin{psmallmatrix}4&3\\6&3\end{psmallmatrix}$, write $A = LU$
with $L = \begin{psmallmatrix}\ell_{11}&0\\\ell_{21}&\ell_{22}
\end{psmallmatrix}$, $U = \begin{psmallmatrix}u_{11}&u_{12}\\0&u_{22}
\end{psmallmatrix}$. Fixing $\ell_{11}=\ell_{22}=1$ (Doolittle's
convention) and expanding $LU$ entrywise gives four equations:
\[
  \ell_{11}u_{11} = 4,\quad \ell_{11}u_{12} = 3,\quad
  \ell_{21}u_{11} = 6,\quad \ell_{21}u_{12} + \ell_{22}u_{22} = 3.
\]
Solving in order: $u_{11} = 4$, $u_{12} = 3$, $\ell_{21} = 6/4 = 1.5$, and
$\ell_{22}u_{22} = 3 - 1.5(3) = -1.5 \Rightarrow u_{22} = -1.5$. So
\[
  \begin{psmallmatrix}4&3\\6&3\end{psmallmatrix} =
  \begin{psmallmatrix}1&0\\1.5&1\end{psmallmatrix}
  \begin{psmallmatrix}4&3\\0&-1.5\end{psmallmatrix}.
\]
\end{example}

\begin{editorialnote}
The original prerequisite slide this example is drawn from has a digit
typo at exactly this step -- it states ``$u_{11} = 14$'' where it means
$u_{11} = 4$ (everything downstream in the original slide is computed
correctly using $u_{11}=4$, so it is an isolated typo, not a repeated
error).
\end{editorialnote}

\subsection{Cholesky Decomposition}
For a Hermitian positive-definite matrix $A$, the \emph{Cholesky
decomposition} is $A = LL^\top$ with $L$ lower triangular and positive
diagonal entries:
\[
  L_{jj} = \sqrt{A_{jj} - \sum_{k=1}^{j-1} L_{jk}^2}, \qquad
  L_{ij} = \frac{1}{L_{jj}}\Big(A_{ij} - \sum_{k=1}^{j-1} L_{ik}L_{jk}\Big)
  \ \text{for } i > j.
\]
Because it needs only one triangular factor (not two), Cholesky is the
standard way to solve $\Sigma x = b$ when $\Sigma$ is a covariance matrix
(always symmetric positive-(semi)definite) -- about twice as fast as a
generic LU solve.

\subsection{QR Decomposition}
$A = QR$, with $Q$ orthogonal ($Q^\top Q = I$) and $R$ upper triangular.
One construction method is Gram-Schmidt: given columns $a_1,\dots,a_n$ of
$A$, set $u_1 = a_1$, $e_1 = u_1/\lVert u_1\rVert$, and iteratively
\[
  u_{k+1} = a_{k+1} - \sum_{j=1}^{k}(a_{k+1}\cdot e_j)e_j, \qquad
  e_{k+1} = \frac{u_{k+1}}{\lVert u_{k+1}\rVert}.
\]
Then $Q = [e_1|\cdots|e_n]$ and $R_{ij} = a_j \cdot e_i$ for $i \le j$
(0 below the diagonal).

\subsection{Singular Value Decomposition (SVD)}
The SVD generalizes eigendecomposition to \emph{any} $n\times p$ matrix
(not just square ones):
\[
  A_{n\times p} = U_{n\times n}\, S_{n\times p}\, V_{p\times p}^\top,
\]
with $U^\top U = I_n$, $V^\top V = I_p$ (both orthogonal). The columns of
$V$ are the eigenvectors of $A^\top A$; the columns of $U$ are the
eigenvectors of $AA^\top$; the singular values in $S$ are the square
roots of the (shared, nonnegative) eigenvalues of $A^\top A$ or $AA^\top$.
SVD is what PCA (Part II, Chapter 5) actually computes under the hood.

\section{Matrix Differentiation}

\begin{definition}[Jacobian]
Let $y = \psi(x)$ with $y$ an $m$-vector and $x$ an $n$-vector. The
\emph{Jacobian} $\dfrac{\partial y}{\partial x}$ is the $m\times n$ matrix
with $(i,j)$ entry $\dfrac{\partial y_i}{\partial x_j}$.
\end{definition}

A handful of identities cover almost every derivative you need in this
course (all assuming $A$ does not depend on the variable being
differentiated with respect to, unless stated otherwise):
\begin{align*}
  y = Ax &\ \Rightarrow\ \frac{\partial y}{\partial x} = A. \\
  \alpha = y^\top A x &\ \Rightarrow\ \frac{\partial \alpha}{\partial x}
    = y^\top A, \qquad \frac{\partial \alpha}{\partial y} = x^\top A^\top. \\
  \alpha = x^\top A x \ (A\text{ not symmetric}) &\ \Rightarrow\
    \frac{\partial \alpha}{\partial x} = x^\top (A + A^\top). \\
  \alpha = x^\top A x \ (A\text{ symmetric}) &\ \Rightarrow\
    \frac{\partial \alpha}{\partial x} = 2x^\top A. \\
  \alpha = x^\top x &\ \Rightarrow\ \frac{\partial \alpha}{\partial x} = 2x^\top.
\end{align*}
The symmetric quadratic-form identity, $\partial(x^\top A x)/\partial x =
2x^\top A$, is the one you will use over and over: it is exactly the
gradient of the ridge-regression penalty $\lambda \beta^\top \beta$ (Part
II, Chapter 3) and of the Gaussian log-density's quadratic term.

\begin{theorem}[Derivative of a matrix inverse]
Let $A(\alpha)$ be a nonsingular $m\times m$ matrix whose entries are
differentiable functions of a scalar $\alpha$. Then
\[
  \frac{\partial A^{-1}}{\partial \alpha}
  = -A^{-1} \frac{\partial A}{\partial \alpha} A^{-1}.
\]
\end{theorem}

\begin{editorialnote}
The original prerequisite slide states this identity \emph{without} the
leading minus sign (as $A^{-1}\frac{\partial A}{\partial\alpha}A^{-1}$).
That is incorrect -- differentiate $AA^{-1}=I$ with respect to $\alpha$
using the product rule to see why the minus sign is required:
$\frac{\partial A}{\partial\alpha}A^{-1} + A\frac{\partial
A^{-1}}{\partial\alpha} = 0 \Rightarrow \frac{\partial
A^{-1}}{\partial\alpha} = -A^{-1}\frac{\partial A}{\partial\alpha}A^{-1}$.
\end{editorialnote}

\section{Summary}
\begin{summarybox}
\begin{itemize}[nosep]
  \item Matrix multiplication requires matching inner dimensions; $(AB)^\top
        = B^\top A^\top$ and $(AB)^{-1} = B^{-1}A^{-1}$ both reverse order.
  \item A matrix is invertible iff its determinant is nonzero; for $2\times2$,
        $A^{-1} = \frac{1}{\det A}\begin{psmallmatrix}d&-b\\-c&a\end{psmallmatrix}$.
  \item Eigenvalues/eigenvectors satisfy $Ax=\lambda x$; for $2\times2$,
        eigenvalues sum to the trace and multiply to the determinant --
        a fast correctness check for eigenvalue code.
  \item LU, Cholesky, and QR decompose a matrix for solving linear systems;
        SVD generalizes eigendecomposition to non-square matrices and
        underlies PCA.
  \item $\partial(x^\top A x)/\partial x = 2x^\top A$ for symmetric $A$ is
        the single most-reused matrix-calculus identity in this course.
\end{itemize}
\end{summarybox}

\section{Exercises}
Three exercises per language (\texttt{exercises/\{python,r,rust\}/}):
\begin{enumerate}
  \item \textbf{Exercise 1 -- Core operations:} matrix multiply, transpose,
        $2\times2$ determinant and inverse, verifying $(AB)^\top=B^\top A^\top$
        and $\det(AB)=\det(A)\det(B)$.
  \item \textbf{Exercise 2 -- LU decomposition:} Doolittle's method for a
        $2\times2$ matrix, verified by reconstructing $A$ from $L$ and $U$.
  \item \textbf{Exercise 3 -- Eigenvalues/eigenvectors:} the $2\times2$
        characteristic-polynomial formula, verified against the
        trace/determinant identities above.
\end{enumerate}
Each has a broken counterpart to debug in
\texttt{exerciseNN\_buggy.\{py,R,rs\}}.

\section*{References}
\begin{itemize}[nosep]
  \item Barnes, R. J. (2009). \emph{Matrix Differentiation (and some
        other stuff)}.
  \item Wikipedia contributors. \emph{Matrix decomposition}. Wikipedia,
        The Free Encyclopedia.
\end{itemize}

\end{document}
